% year: תש"פ
% type: מבחן
% semester: א
% moed: א
% part: א
% number: 1
% points: 20
% categories: func-analysis, derivatives, functions
% source: exams/מבחנים/תשפ/מועד א תשפ.pdf

\begin{question}
חקרו את הפונקציה $y=\frac{e^{2x}}{x}$ לפי הסעיפים הבאים:
\begin{enumerate}
  \item תחום הגדרה, נקודות חיתוך עם הצירים, זוגיות, אי-זוגיות
  \item רציפות, נקודות אי-רציפות
  \item אסימפטוטות
  \item תחומי עליה וירידה, נקודות קיצון
  \item תחומי קמירות וקעירות, נקודות פיתול
  \item תיאור גרפי.
\end{enumerate}
\end{question}

\begin{hint}
$y'=\frac{e^{2x}(2x-1)}{x^2}$ ו-$y''=\frac{2e^{2x}(2x^2-2x+1)}{x^3}$. בדקו את הסימן של כל גורם.
\end{hint}

\begin{solution}
\begin{enumerate}
  \item \textbf{תחום הגדרה:} $x\neq0$, כלומר $(-\infty,0)\cup(0,\infty)$.

  \textbf{חיתוך עם הצירים:} $x=0$ לא בתחום — אין חיתוך עם ציר $y$. $e^{2x}>0$ תמיד, ולכן $y\ne0$ — אין חיתוך עם ציר $x$. (הסימן: $y<0$ עבור $x<0$ ו-$y>0$ עבור $x>0$.)

  \textbf{זוגיות:} $y(-x)=\frac{e^{-2x}}{-x}$. למשל $y(1)=e^2$, $y(-1)=-e^{-2}$, ולכן $y(-1)\ne y(1)$ וגם $y(-1)\ne -y(1)$. הפונקציה \textbf{אינה זוגית ואינה אי-זוגית}.

  \item הפונקציה היא מנה של פונקציות רציפות ($e^{2x}$ ו-$x$) ולכן רציפה בכל נקודה שבה המכנה שונה מאפס, כלומר בכל תחום הגדרתה. הנקודה היחידה שבעייתית היא $x=0$ (שבה הפונקציה אינה מוגדרת):
  \[ \lim_{x\to0^-}\frac{e^{2x}}{x}=\frac{1}{0^-}=-\infty,\qquad \lim_{x\to0^+}\frac{e^{2x}}{x}=\frac{1}{0^+}=+\infty, \]
  ולכן $x=0$ היא נקודת אי-רציפות מסוג שני (עיקרית).

  \item \textbf{אסימפטוטה אנכית:} $x=0$ (לפי הגבולות בסעיף הקודם).

  \textbf{ב-$-\infty$:} $e^{2x}\to0$ ו-$\frac1x\to0$, לכן $\lim_{x\to-\infty}y=0$ — אסימפטוטה אופקית $y=0$ משמאל (הגרף מתחת לציר).

  \textbf{ב-$+\infty$:} $\frac{y}{x}=\frac{e^{2x}}{x^2}\to\infty$ (לפי לופיטל פעמיים: $\lim\frac{e^{2x}}{x^2}=\lim\frac{2e^{2x}}{2x}=\lim\frac{4e^{2x}}{2}=\infty$), לכן אין אסימפטוטה אופקית או משופעת מימין.

  \item לפי כלל המנה:
  \[ y'=\frac{2e^{2x}\cdot x - e^{2x}}{x^2} = \frac{e^{2x}(2x-1)}{x^2}. \]
  מכיוון ש-$e^{2x}>0$ ו-$x^2>0$, הסימן של $y'$ הוא הסימן של $2x-1$:
  \begin{itemize}
    \item $y'<0$ עבור $x<0$ ועבור $0<x<\frac12$: הפונקציה \textbf{יורדת} ב-$(-\infty,0)$ וב-$(0,\frac12)$.
    \item $y'>0$ עבור $x>\frac12$: הפונקציה \textbf{עולה} ב-$(\frac12,\infty)$.
  \end{itemize}
  ב-$x=\frac12$ הנגזרת מחליפה סימן מ-$-$ ל-$+$, ולכן זו נקודת \textbf{מינימום מקומי}: $y(\tfrac12)=\frac{e}{1/2}=2e$, כלומר $(\frac12,2e)$. אין נקודות קיצון נוספות.

  \item נגזור את $y'=e^{2x}\cdot\frac{2x-1}{x^2}$:
  \[ \left(\frac{2x-1}{x^2}\right)'=\frac{2x^2-2x(2x-1)}{x^4}=\frac{2-2x}{x^3}, \]
  \[ y''=e^{2x}\left(\frac{2(2x-1)}{x^2}+\frac{2-2x}{x^3}\right)=\frac{e^{2x}\left(4x^2-2x+2-2x\right)}{x^3}=\frac{2e^{2x}\left(2x^2-2x+1\right)}{x^3}. \]
  לטרינום $2x^2-2x+1$ דיסקרימיננטה $4-8<0$ ומקדם מוביל חיובי, ולכן הוא חיובי תמיד. מכאן הסימן של $y''$ הוא הסימן של $x^3$, כלומר של $x$:
  \begin{itemize}
    \item $x<0$: $y''<0$ — הפונקציה \textbf{קעורה} ($\cap$) ב-$(-\infty,0)$.
    \item $x>0$: $y''>0$ — הפונקציה \textbf{קמורה} ($\cup$) ב-$(0,\infty)$.
  \end{itemize}
  סימן $y''$ מתחלף רק ב-$x=0$, שאינה בתחום ההגדרה, ולכן \textbf{אין נקודות פיתול}.

  \item \textbf{תיאור גרפי:} משמאל לציר $y$ הגרף מתחת לציר $x$: מתחיל קרוב ל-$y=0$ ב-$-\infty$, יורד (וקעור) ושואף ל-$-\infty$ כאשר $x\to0^-$. מימין לציר $y$ הגרף מעל ציר $x$: יורד מ-$+\infty$ (ליד $x=0^+$) עד נקודת המינימום $(\frac12,2e)\approx(0.5,\,5.44)$, ואז עולה (בקמירות) ל-$+\infty$ במהירות אקספוננציאלית.
\end{enumerate}
\end{solution}
